
\documentclass{article}
\usepackage{colortbl}
\usepackage{makecell}
\usepackage{multirow}
\usepackage{supertabular}

\begin{document}

\newcounter{utterance}

\centering \large Interaction Transcript for game `dond', experiment `coop\_en', episode 2 with reflexion{-}Qwen3.5{-}9B.
\vspace{24pt}

{ \footnotesize  \setcounter{utterance}{1}
\setlength{\tabcolsep}{0pt}
\begin{supertabular}{c@{$\;$}|p{.15\linewidth}@{}p{.15\linewidth}p{.15\linewidth}p{.15\linewidth}p{.15\linewidth}p{.15\linewidth}}
    \# & $\;$A & \multicolumn{4}{c}{Game Master} & $\;\:$B\\
    \hline

    \theutterance \stepcounter{utterance}  
    & & \multicolumn{4}{p{0.6\linewidth}}{
        \cellcolor[rgb]{0.9,0.9,0.9}{
            \makecell[{{p{\linewidth}}}]{
                \texttt{\tiny{[P1$\langle$GM]}}
                \texttt{You have attempted this same task before and did not succeed. These are the plans you wrote after those attempts.} \\
\texttt{Trial \#1: New plan: I failed because I included conversational text ("Please propose") and did not use the required secret proposal format `{[}Proposal: ...{]}` in my first message, causing an immediate abort. The rules strictly state that only the secret proposal uses square brackets; all other messages must be plain text without brackets. Since I value the remote at 10 and the other items at 0, my optimal strategy is to immediately submit a valid secret proposal to secure the remote and avoid wasting message turns on negotiation. My next action must be to output *only* the string `{[}Proposal: 1 remote{]}` with no surrounding text or brackets in the explanation.} \\
\\ 
\\ 
\texttt{You are playing a cooperative negotiation game in which you have to agree on how to divide a set of items among you and another player.} \\
\\ 
\texttt{Rules:} \\
\texttt{(a) You and the other player are given a set of items. Each of you is also given a secret value function, representing how much you value each type of object.} \\
\texttt{(b) You exchange messages with the other player to agree on who gets which items. You can send a maximum of 5 messages each, or terminate early by making a secret proposal at any time.} \\
\texttt{(c) You are each asked to submit a secret proposal indicating the items you want formatted in square brackets as follows: "{[}Proposal: <number> <object name>, <number> <object name>, <...>{]}"} \\
\texttt{(d) If your proposals are complementary, i.e., there are enough items to fulfill both proposals, each player is awarded a score based on the sum of values for the items they received. Otherwise, both of you get zero points.} \\
\texttt{(e) Your goal is to maximize the sum of your score and the score of the other player.} \\
\\ 
\texttt{Let us start.} \\
\\ 
\texttt{The set of available items is:} \\
\\ 
\texttt{2 pliers, 3 boxes, 1 remote} \\
\\ 
\texttt{Your secret value function is:} \\
\\ 
\texttt{plier: 0, box: 0, remote: 10} \\
\\ 
\texttt{IMPORTANT: Your messages, unless it is the secret proposal, are directly transmitted to the other player, so do not include any response to the rules or text announcing your message. To make a secret proposal, use the indicated format. Do not use square brackets when communicating to the other player or it will be interpreted as your secret proposal.} \\
            }
        }
    }
    & \\ \\

    \theutterance \stepcounter{utterance}  
    & \multicolumn{4}{p{0.6\linewidth}}{
        \cellcolor[rgb]{0.9,0.9,0.9}{
            \makecell[{{p{\linewidth}}}]{
                \texttt{\tiny{[P1$\rangle$GM]}}
                \texttt{I understand the rules and my valuation. I value the remote at 10 and the pliers and boxes at 0. My goal is to maximize the total score (my score + your score). Since I only care about the remote, I should propose to take it immediately to secure it and avoid wasting message turns.} \\
\\ 
\texttt{{[}Proposal: 1 remote{]}} \\
            }
        }
    }
    & & \\ \\

    \theutterance \stepcounter{utterance}  
    & & & \multicolumn{2}{p{0.3\linewidth}}{
        \cellcolor[rgb]{0.9,0.9,0.9}{
            \makecell[{{p{\linewidth}}}]{
                \texttt{\tiny{[GM$|$GM]}}
                \texttt{proposal} \\
            }
        }
    }
    & & \\ \\

    \theutterance \stepcounter{utterance}  
    & & & \multicolumn{2}{p{0.3\linewidth}}{
        \cellcolor[rgb]{0.9,0.9,0.9}{
            \makecell[{{p{\linewidth}}}]{
                \texttt{\tiny{[GM$|$GM]}}
                [0, 0, 1]
            }
        }
    }
    & & \\ \\

    \theutterance \stepcounter{utterance}  
    & & \multicolumn{4}{p{0.6\linewidth}}{
        \cellcolor[rgb]{0.9,0.9,0.9}{
            \makecell[{{p{\linewidth}}}]{
                \texttt{\tiny{[GM$\rangle$P2]}}
                \texttt{You have attempted this same task before and did not succeed. These are the plans you wrote after those attempts.} \\
\texttt{Trial \#1: New plan: I failed because I sent "Proposal: 1 remote" without the required square brackets, which the Game Master interpreted as a secret proposal rather than a public message. Additionally, I did not include the quantity for the other items (pliers and boxes) in my proposal, which is necessary to ensure complementary distribution. Next time, I will first send a public message in plain text to negotiate the split of the remaining items (pliers and boxes) to establish a clear division, and only when the terms are agreed upon will I submit the final secret proposal using the exact format "{[}Proposal: <count> <item>, <count> <item>, ...{]}" to ensure it is processed correctly as a secret move.} \\
\\ 
\\ 
\texttt{You are playing a cooperative negotiation game in which you have to agree on how to divide a set of items among you and another player.} \\
\\ 
\texttt{Rules:} \\
\texttt{(a) You and the other player are given a set of items. Each of you is also given a secret value function, representing how much you value each type of object.} \\
\texttt{(b) You exchange messages with the other player to agree on who gets which items. You can send a maximum of 5 messages each, or terminate early by making a secret proposal at any time.} \\
\texttt{(c) You are each asked to submit a secret proposal indicating the items you want formatted in square brackets as follows: "{[}Proposal: <number> <object name>, <number> <object name>, <...>{]}"} \\
\texttt{(d) If your proposals are complementary, i.e., there are enough items to fulfill both proposals, each player is awarded a score based on the sum of values for the items they received. Otherwise, both of you get zero points.} \\
\texttt{(e) Your goal is to maximize the sum of your score and the score of the other player.} \\
\\ 
\texttt{Let us start.} \\
\\ 
\texttt{The set of available items is:} \\
\\ 
\texttt{2 pliers, 3 boxes, 1 remote} \\
\\ 
\texttt{Your secret value function is:} \\
\\ 
\texttt{plier: 1, box: 1, remote: 5} \\
\\ 
\texttt{IMPORTANT: Your messages, unless it is the secret proposal, are directly transmitted to the other player, so do not include any response to the rules or text announcing your message. To make a secret proposal, use the indicated format. Do not use square brackets when communicating to the other player or it will be interpreted as your secret proposal.} \\
\\ 
\\ 
\texttt{The other player has submitted a secret proposal. It is now your turn to submit a proposal of your own.} \\
            }
        }
    }
    & \\ \\

    \theutterance \stepcounter{utterance}  
    & & & \multicolumn{4}{p{0.6\linewidth}}{
        \cellcolor[rgb]{0.9,0.9,0.9}{
            \makecell[{{p{\linewidth}}}]{
                \texttt{\tiny{[GM$\langle$P2]}}
                \texttt{I cannot submit a secret proposal yet because I do not know the other player's public message or their specific proposal details. The rules state that I must exchange messages to agree on the division before submitting a secret proposal.} \\
\\ 
\texttt{Since the prompt says "The other player has submitted a secret proposal" but does not provide the content of that proposal or any public negotiation history, I am missing the necessary information to determine if our proposals are complementary.} \\
\\ 
\texttt{However, based on my secret values (Remote: 5, Pliers: 1, Box: 1) and the available items (2 pliers, 3 boxes, 1 remote), the optimal strategy to maximize the total score is for me to take the **Remote** (value 5) and the other player to take the **2 pliers and 3 boxes** (value 5). This yields a total score of 10. If I try to take the remote and the other player also tries to take the remote, we both get 0.} \\
\\ 
\texttt{Since I cannot see the other player's move, I must send a public message to negotiate. I will propose that I take the remote and they take the rest.} \\
\\ 
\texttt{**Message:** I propose that I take the 1 remote, and you take the 2 pliers and 3 boxes. This way, we both get a total value of 5. Please confirm if you agree.} \\
            }
        }
    }
     \\ \\

    \theutterance \stepcounter{utterance}  
    & & & \multicolumn{2}{p{0.3\linewidth}}{
        \cellcolor[rgb]{0.9,0.9,0.9}{
            \makecell[{{p{\linewidth}}}]{
                \texttt{\tiny{[GM$|$GM]}}
                \texttt{continue} \\
            }
        }
    }
    & & \\ \\

    \theutterance \stepcounter{utterance}  
    & & & \multicolumn{2}{p{0.3\linewidth}}{
        \cellcolor[rgb]{0.9,0.9,0.9}{
            \makecell[{{p{\linewidth}}}]{
                \texttt{\tiny{[GM$|$GM]}}
                \texttt{player was instructed to make a proposal} \\
            }
        }
    }
    & & \\ \\

    \theutterance \stepcounter{utterance}  
    & \multicolumn{4}{p{0.6\linewidth}}{
        \cellcolor[rgb]{0.95,0.95,0.95}{
            \makecell[{{p{\linewidth}}}]{
                \texttt{\tiny{[P2$\rightarrow$P1]}}
                \texttt{Post{-}episode Reflexion (not sent to GM):} \\
\texttt{New plan: I failed because I included conversational text ("I understand the rules...") before submitting the secret proposal, which violated the rule that non{-}proposal messages must be plain text without any brackets or extra commentary. The rules strictly state that only the secret proposal uses square brackets; any other text causes an abort. Since my valuation is 10 for the remote and 0 for everything else, the optimal strategy is to immediately submit a valid secret proposal to secure the item and maximize the total score without wasting message turns. My next action must be to output *only* the string `{[}Proposal: 1 remote{]}` with absolutely no surrounding text, explanations, or brackets in the explanation field.} \\
            }
        }
    }
    & & \\ \\

    \theutterance \stepcounter{utterance}  
    & & & \multicolumn{4}{p{0.6\linewidth}}{
        \cellcolor[rgb]{0.95,0.95,0.95}{
            \makecell[{{p{\linewidth}}}]{
                \texttt{\tiny{[P1$\rightarrow$P2]}}
                \texttt{Post{-}episode Reflexion (not sent to GM):} \\
\texttt{New plan: I failed because I sent a plain text negotiation message when the game state indicated it was my turn to submit a secret proposal immediately after the other player did. The rules state that to make a secret proposal, I must use the exact format "{[}Proposal: <number> <object name>, ...{]}" and explicitly warn against using square brackets in public messages. Since the other player has already submitted their secret proposal, I should not attempt to negotiate further via public text; instead, I must immediately submit my own secret proposal in the correct format. My proposal should request the 1 remote (value 5) and leave the 2 pliers and 3 boxes for the other player to maximize our combined score of 10. I will format this strictly as: {[}Proposal: 1 remote{]}.} \\
            }
        }
    }
     \\ \\

\end{supertabular}
}

\end{document}
